% Generated draft: requires mathematical and source review.
\documentclass[11pt]{article}
\usepackage[margin=25mm]{geometry}
\usepackage{fontspec}
\setmainfont{Latin Modern Roman}
\usepackage{amsmath,amssymb,mathtools,unicode-math}
\setmathfont{Latin Modern Math}
\usepackage{enumitem,needspace,xurl,hyperref}
\setlength{\parindent}{0pt}
\setlength{\parskip}{5pt}
\setlength{\emergencystretch}{3em}
\newcommand{\N}{\mathbb{N}}
\newcommand{\Z}{\mathbb{Z}}
\newcommand{\Q}{\mathbb{Q}}
\newcommand{\R}{\mathbb{R}}
\newcommand{\C}{\mathbb{C}}
\newcommand{\F}{\mathbb{F}}
\newcommand{\A}{\mathbb{A}}
\newcommand{\PP}{\mathbb{P}}
\newcommand{\E}{\mathbb{E}}
\newcommand{\eps}{\varepsilon}
\newcommand{\dd}{\,\mathrm{d}}
\begin{document}
\section*{UnsolvedMath ID: 7500114}
\textbf{Source problem code:} catalogue entry
\subsection{Definitions and mathematical statement}
For a countable partially ordered set \((P,\preceq)\), use the convention that a filter is a nonempty upward-closed subset \(F\subseteq P\) which is downward directed: for \(p,q\in F\) there is \(r\in F\) with \(r\preceq p,q\). Let \(\operatorname{MF}(P)\) be the set of maximal filters, with basic open sets
\[
 U_p=\{F\in\operatorname{MF}(P):p\in F\}\qquad(p\in P).
\]
This is the usual maximal-filter-space interpretation of ``MF space''. A countable basis does not mean the set of points is countable.

Classically, for a noncompact locally compact Hausdorff space \(X\), the Alexandroff one-point compactification adjoins a point \(\infty\); it keeps the open sets of \(X\), and neighborhoods of \(\infty\) have the form
\[
 \{\infty\}\cup(X\setminus K),\qquad K\subset X\text{ compact and closed}.
\]
The result is compact Hausdorff and contains \(X\) as a dense open subspace. The source asks for the reverse-mathematical strength of the corresponding theorem for countably based MF spaces.

A precise reverse-mathematical target must fix the base subsystem, the codes for posets, spaces, compactness and maps, and whether the compactification must itself be an MF space with a supplied countable presentation. It must also state the local compactness/Hausdorff assumptions and treatment of already compact spaces. The source excerpt includes none of these; the classical hypotheses above are explanatory, not silently inserted as the exact formal principle. Over a conventional base such as \(\mathrm{RCA}_0\), the intended task would be to identify an equivalent axiom system or otherwise locate the exact strength of that fully coded principle, but the base is not expressly given.

The supplied mid-2026 assessment says no verified solution was found, and the 2026-07-29 nested-list audit says the surviving formulation appears open and needs review. These are dated source reports, not an independently verified classification of the formal principle.
\subsection{Short English statement}
Identify the axioms needed to prove one-point compactification for the intended class of countably based maximal-filter spaces. The precise topological hypotheses, output representation and reverse-mathematical coding/base theory must be recovered from the original question.
\subsection{Sources}
\begin{itemize}
\item Local registry: \url{data/unsolvedmath/problems.json}, permanent UnsolvedMath ID \(7500114\), source code catalogue entry; supplied statement, background, and research summary, updated 2026-08-06.
\item Montalbán, \emph{Open Questions in Reverse Mathematics} (2011), Question 14: \url{https://math.berkeley.edu/~antonio/papers/questionsRM.pdf}. Supplied author-preprint PDF text plus manual normalization, accessed 2026-07-29.
\item The supplied nested-list audit records ``appears open'' and continuing review need; no separate resolution theorem or full coding specification is supplied.
\end{itemize}
\end{document}
